% ============================================================
%  Notae (быстрый старт)
%
%  Это обычный .tex-файл. Всё, что вне {...} и не начинается
%  со служебного слова, считается комментарием/мусором.
%  Полное описание — в README.md.
%
%  ИМЕНА СЛАЙДОВ (нумерация с единицы):
%    Title    — титульный слайд (если заполнен);
%    Slide01  — первая сцена;
%    Slide02  — вторая и т.д.
%  Эти же имена видны в терминале и в файлах media/slides/.
% ============================================================

\documentclass[a4paper,11pt]{article}

% --- кириллица и шрифты ---
\usepackage{cmap}
\usepackage[T2A]{fontenc}
\usepackage[utf8]{inputenc}
\usepackage[russian]{babel}

\usepackage{amsmath, amsfonts, amssymb, amsthm}
\usepackage{mathrsfs}
\usepackage{geometry}
\usepackage{xcolor}

% ------------------------------------------------------------
%  СВОИ МАКРОСЫ (необязательно)
%
%  Notae переносит их в генерируемую преамбулу, и они работают
%  во всех строках. Поддерживаются:
%    \newcommand, \renewcommand, \providecommand,
%    \DeclareMathOperator, \newtheorem,
%    \newenvironment, \renewenvironment, \let.
% ------------------------------------------------------------
\newcommand{\R}{\mathbb{R}}
\newcommand{\N}{\mathbb{N}}
\newcommand{\norm}[1]{\left\lVert #1 \right\rVert}
\DeclareMathOperator{\E}{E}
\DeclareMathOperator*{\argmin}{arg\,min}

\newtheorem{theorem}{Теорема}
\newtheorem{lemma}{Лемма}

\newenvironment{note}{\begin{quote}\itshape}{\end{quote}}

% Ещё примеры (раскомментируйте при необходимости):
% \newcommand{\D}{\mathcal{D}}
% \let\oldphi\phi
% \usepackage{tikz}

\begin{document}

% ============================================================
%  ТИТУЛЬНЫЙ СЛАЙД (необязательно)
%
%  Любое поле можно оставить пустым {} — оно не выводится.
%  pause: — сколько секунд держать титульник.
%  Если заполнено хотя бы одно поле, слайд называется Title.
% ============================================================
title: {Название доклада}
author: {ФИО}
affiliation: {Институт}
event: {Название конференции}
date: {Дата}
pause: 5.0


% ============================================================
%  ОБЩИЙ БЛОК (необязательно)
%
%  Объявления до первой scene: доступны во всех сценах:
%  общие строки, субтитры и стили.
% ============================================================
common:
line gap: {\vspace{1.0cm}}

% Стиль строк по умолчанию:
style default: size=36 color=TEXT anim=write align=center
% Стиль субтитров по умолчанию:
style sub default: size=26 color=MUTED


% ============================================================
%  КРАТКАЯ СПРАВКА
%
%  T: <задержка> <имя>        показать строку или субтитр
%  T: <задержка> <имя> -      убрать строку
%  T: <задержка> sub: -       убрать субтитр
%  pause: <секунды>           просто подождать
%
%  Объявление строки прямо в таймлайне (не рекомендуется):
%  T: 0.5 line lx: {Строка, объявленная на месте.}
%
%  ТЕКСТ И ФОРМУЛЫ. Строка с обычными словами (кириллица,
%  английский, диакритика) — текст. Чистая математика
%  оборачивается в формулу автоматически; надёжнее писать
%  её явно через $...$ или \[...\].
%
%  ДЛИНА СТРОК. Поля презентации — по 2 см слева и справа;
%  слишком длинная строка переносится автоматически, но лучше
%  держать строки короткими. Субтитр — не более 3 строк.
%
%  ПРОВЕРКА. Notae различает ошибки (из-за них слайд может
%  не собраться) и предупреждения (длинные строки, субтитры,
%  плотность); предупреждения сборку не срывают.
%
%  ВАЖНО. Служебное слово в начале строки обрывает {...}
%  только тогда, когда перед ним уже стоит закрывающая }.
%  Пробелы, пустые строки и необязательный \\ после } не
%  мешают. Поэтому строки, похожие на команды, можно писать
%  внутри ещё не закрытого {...}. Если содержимое не закрыто,
%  а дальше идёт настоящее служебное слово (T:, line:, scene:),
%  скрипт подскажет, где поставить }.
%  Оформление }\\ перед служебным словом разрешено.
%  Если содержимое само кончается на }, поставьте две скобки: ...}}.
% ============================================================

% ============================================================
%                     Scene01 – ШАБЛОН СЦЕНЫ 
% ============================================================

% Объявление сцены
scene: {}

% Описание содержимого
line l1: {}

% Объявление субтитров
sub s1: {Субтитр 1.}

% Таймлайн
T: 1.5 l1

T: 1.0 s1

pause: 5.0




\end{document}
